\section{Results}
\label{sec:results}
\subsection{Classification Is Sensitive to Fold and Endpoint}
Simple baselines reveal substantial fold dependence. The colour probe obtains macro-F1 values of 0.88696, 0.46405, and 0.12322 on folds 0, 1, and 2. The histogram-of-oriented-gradients support-vector machine obtains 0.79602, 0.18182, and 0.98412. These patterns motivate scrutiny of capture-specific cues, but do not identify a unique shortcut mechanism. A historical random-initialised ResNet-18 trained for 15 epochs averages 0.82135, whereas the later matched nine-run, 60-epoch ResNet-50 control averages 0.82332 at the final endpoint. The two controls differ in backbone and budget and are not interchangeable.

Table~\ref{tab:architectures} and Fig.~\ref{fig:endpoints} show the retained sweep. MobileNetV4 has a selected mean macro-F1 of 0.99746 and final mean of 0.91310. The endpoint gap is larger for EfficientNetV2-S (0.97615 versus 0.70874) and VGG16-BN (0.89194 versus 0.45870), while ConvNeXtV2-Tiny is comparatively stable (0.90891 versus 0.90238). Consequently, a near-perfect selected score cannot stand alone as a performance summary. These values include the historically flagged folds and do not establish external generalisation, physical-wear recognition, or the optimal checkpoint policy for new tyres.

\begin{table}[t]
\caption{Archival classification summary in alphabetical order. Each row averages nine runs (three grouped folds and three seeds). These internal validation summaries are retained for completeness; fold-specific analysis is primary.}
\label{tab:architectures}
\centering
\small
\setlength{\tabcolsep}{4pt}
\begin{tabular}{lrr}
\toprule
Architecture & Selected & Final \\
\midrule
CLIP ViT-B/16 & 0.95076 & 0.78897 \\
CoAtNet-0 & 0.80537 & 0.72041 \\
ConvNeXtV2-T & 0.90891 & 0.90238 \\
DeiT-III-S & 0.74726 & 0.72855 \\
DenseNet-121 & 0.89660 & 0.83852 \\
DINOv2-B & 0.91277 & 0.88380 \\
DINOv2-S & 0.87346 & 0.78842 \\
EfficientNetV2-S & 0.97615 & 0.70874 \\
MaxViT-T & 0.92189 & 0.85982 \\
MobileNetV4 & 0.99746 & 0.91310 \\
RegNetY-016 & 0.93495 & 0.74087 \\
ResNet-50 & 0.93651 & 0.82922 \\
ResNeXt-50 & 0.92260 & 0.76566 \\
Swin-S & 0.95259 & 0.85744 \\
Swin-T & 0.97673 & 0.83478 \\
VGG16-BN & 0.89194 & 0.45870 \\
ViT-S & 0.93235 & 0.75858 \\
\bottomrule
\end{tabular}
\end{table}

\begin{figure*}[t]
\centering
\includegraphics[width=.88\textwidth]{endpoints.pdf}
\caption{Validation-selected and fixed-final macro-F1 means for all 17 retained configurations. Each point averages nine runs on the same historical folds; connecting segments highlight endpoint differences and are not uncertainty intervals. No test-set architecture ranking is implied.}
\label{fig:endpoints}
\end{figure*}

\subsection{Explanation Diagnostics and Same-Fold Confirmation}
The inherited hypothesis that normalised tread evidence predicts cross-fold stability better than accuracy is not supported: the recorded primary and reference statistics are 0.42710 and 0.90743, respectively, across ten entries. This analysis uses legacy selected-endpoint statistics and is not relabelled as a fixed-budget result. The marking/damage hypothesis is inconclusive or undefined, and the fine-grained architecture hypothesis remains untested because its planned model arms are absent. In particular, fold-1 marking removal operates on zero active marking images; its near-zero effect cannot establish marking independence.

Table~\ref{tab:confirmation} reports the confirmation stage. Random initialisation is negative for both architectures at both endpoints. Class-weighted sampling has a positive selected effect for ConvNeXtV2-Tiny but a negative selected effect for MobileNetV4, while both final effects are positive. Uniform sampling is not consistently beneficial. The supported interpretation therefore specifies architecture, endpoint, and the existing fold. The experiment neither isolates factor interactions nor establishes a population-level sampling benefit.
\begin{table}[t]
\caption{Same-fold confirmation: mean macro-F1 change relative to each matched baseline over three seeds. Selection and final endpoints can give different sampling conclusions.}
\label{tab:confirmation}
\centering
\small
\setlength{\tabcolsep}{4pt}
\begin{tabular}{llrr}
\toprule
Architecture & Factor & $\Delta$ selected & $\Delta$ final \\
\midrule
ConvNeXtV2-T & Class-weighted & +0.01257 & +0.02396 \\
ConvNeXtV2-T & Uniform & +0.00000 & +0.01027 \\
ConvNeXtV2-T & Random init. & -0.42857 & -0.38737 \\
MobileNetV4 & Class-weighted & -0.02839 & +0.01886 \\
MobileNetV4 & Uniform & -0.06959 & -0.17079 \\
MobileNetV4 & Random init. & -0.37150 & -0.27525 \\
\bottomrule
\end{tabular}
\end{table}


\subsection{Region Overlap and Boundary Fidelity Differ}
All 81 localisation runs complete 60 epochs and evaluation, corresponding to 4,860 epoch records. Table~\ref{tab:localisation} restricts the principal summary to fold 1; this avoids treating flagged-fold averages as the main estimate but does not make fold 1 an external test. Semantic models have tread IoU between approximately 0.972 and 0.976, whereas their two-pixel boundary F1 is approximately 0.358--0.374. YOLO instance-segmentation models obtain tread IoU near 0.856--0.858 and boundary F1 near 0.189--0.226.

Figure~\ref{fig:localisation} makes the metric distinction explicit. Region overlap rewards agreement across a large interior area; boundary F1 is sensitive to contour displacement under its narrow tolerance. Neither metric alone establishes that the model has recovered the information needed for proxy classification, let alone physical tread condition. The evaluation is against one existing manual reference set rather than an independently repeated annotation standard.
\begin{table*}[t]
\caption{Fold-1 localisation means over three seeds. AP denotes COCO-style AP at IoU thresholds 0.50:0.95. Boundary F1 uses a two-native-pixel tolerance. Dashes denote undefined mask metrics for box-only models.}
\label{tab:localisation}
\centering
\small
\setlength{\tabcolsep}{4pt}
\begin{tabular}{lrrrr}
\toprule
Model & Box AP & Mask AP & Tread IoU & Boundary F1 \\
\midrule
DeepLabV3+/R34 & 0.98612 & 0.99675 & 0.97176 & 0.35988 \\
RT-DETRv2-R18 & 0.94073 & -- & -- & -- \\
SegFormer-B0 & 0.91189 & 0.99836 & 0.97372 & 0.37258 \\
SegFormer-B2 & 0.97657 & 0.99982 & 0.97569 & 0.37412 \\
U-Net/R34 & 0.91413 & 0.99704 & 0.97216 & 0.35820 \\
YOLO26-n det. & 0.95278 & -- & -- & -- \\
YOLO26-n seg. & 0.95204 & 0.95890 & 0.85815 & 0.18926 \\
YOLO26-s det. & 0.96081 & -- & -- & -- \\
YOLO26-s seg. & 0.95227 & 0.96215 & 0.85627 & 0.22623 \\
\bottomrule
\end{tabular}
\end{table*}

\begin{figure*}[t]
\centering
\includegraphics[width=\textwidth]{localisation.pdf}
\caption{Fold-1 semantic/instance-segmentation results: (a) tread IoU and (b) boundary F1 at two native pixels. Points are means and whiskers are sample standard deviations across three training seeds. The panels deliberately use different horizontal scales. Seed variation is not a confidence interval over independent tyres.}
\label{fig:localisation}
\end{figure*}

\subsection{Predicted Crops and Fixed Fusion Have Negative Mean Effects}
The paired frozen-classifier analysis does not show an overall gain from the predicted crops. Mean macro-F1 changes are -0.01257 for tyre-only and -0.01467 for tread-only inputs across the 81 localisation source runs. These means do not imply that every configuration/fold pair worsens. They quantify the declared intervention on a classifier trained for full-image inputs, with no crop-specific retraining.

Fixed probability fusion also has negative mean effects (Table~\ref{tab:fusion}, Fig.~\ref{fig:fusion}). Equal tyre/tread fusion gives -0.01568; equal full/tyre/tread fusion gives -0.01324. The latter is a decrease of about 1.324 macro-F1 percentage points, not a 1.324\% relative error reduction. The result argues against assuming that localisation or additional views automatically help an existing classifier. Lost context, changed scale, crop distribution shift, localisation errors, and correlated predictions are plausible explanations, but the present design does not separate their contributions. Crop-trained recognition and learned fusion would require new experiments.
\begin{table}[t]
\caption{Fixed probability fusion: mean paired macro-F1 changes over 81 source runs. Repeated images and shared classifiers make these rows dependent.}
\label{tab:fusion}
\centering
\small
\setlength{\tabcolsep}{4pt}
\begin{tabular}{lrr}
\toprule
Arm & Rows & Mean $\Delta$ \\
\midrule
Full frame & 81 & +0.00000 \\
Tyre only & 81 & -0.01257 \\
Tread only & 81 & -0.01467 \\
Tyre + tread & 81 & -0.01568 \\
Full + tyre + tread & 81 & -0.01324 \\
\bottomrule
\end{tabular}
\end{table}

\begin{figure}[t]
\centering
\includegraphics[width=\columnwidth]{fusion.pdf}
\caption{Mean paired changes from the full-frame classifier for four fixed crop/fusion arms. Bars average 81 dependent source runs; no independent-sample uncertainty or favourable-arm selection is implied.}
\label{fig:fusion}
\end{figure}

\subsection{Confidence Diagnostics Are Fold-Dependent}
Temperature scaling leaves the recorded macro-F1 unchanged and reduces confidence-error metrics on the retained calibration-test subsets (Table~\ref{tab:calibration}). On fold 1, expected calibration error decreases from 0.18955 to 0.07377, and negative log-likelihood from 0.35630 to 0.16226, with 41 test images. Extremely small calibrated errors on the flagged folds are not evidence of deployment-grade certainty.

At nominal 90\% coverage, the observed conformal coverages are 0.86885, 0.97561, and 0.93939 on folds 0, 1, and 2, respectively (Table~\ref{tab:conformal}). Thus only fold 0 is below nominal in these observed subsets. Mean set sizes below one on folds 0 and 2 imply empty sets. The inherited abstention field does not count all empty-set events and cannot be presented as a complete rejection policy. Small subset sizes, dependent acquisition, and reuse of internal evidence constrain generalisation of these diagnostics.
\begin{table}[t]
\caption{Confidence diagnostics on held-out calibration-test subsets. ECE is the recorded expected calibration error; values near zero on internal folds do not imply external reliability.}
\label{tab:calibration}
\centering
\small
\setlength{\tabcolsep}{4pt}
\begin{tabular}{llrrrr}
\toprule
Fold & Scores & $n$ & F1 & ECE & NLL \\
\midrule
0 & Raw & 61 & 1.0000 & 0.2115 & 0.2425 \\
0 & Temp. & 61 & 1.0000 & 0.0000 & 0.0000 \\
1 & Raw & 41 & 0.8030 & 0.1896 & 0.3563 \\
1 & Temp. & 41 & 0.8030 & 0.0738 & 0.1623 \\
2 & Raw & 33 & 1.0000 & 0.1293 & 0.1438 \\
2 & Temp. & 33 & 1.0000 & 0.0000 & 0.0000 \\
\bottomrule
\end{tabular}
\end{table}

\begin{table}[t]
\caption{Prediction-set diagnostics at nominal 90\% coverage. The recorded abstention field does not count every empty-set event.}
\label{tab:conformal}
\centering
\small
\setlength{\tabcolsep}{4pt}
\begin{tabular}{rrrrrr}
\toprule
Fold & $n_{cal}$ & $n_{test}$ & Coverage & Set size & Abstain \\
\midrule
0 & 62 & 61 & 0.8689 & 0.8689 & 0.0000 \\
1 & 43 & 41 & 0.9756 & 1.0976 & 0.0976 \\
2 & 35 & 33 & 0.9394 & 0.9394 & 0.0000 \\
\bottomrule
\end{tabular}
\end{table}


\subsection{Matched Point Learning Reduces Still-Image Error}
Table~\ref{tab:geometry} reports the fixed-final comparison, independently recomputed here from the 432 paired point records. Mean horizontal error is 1.312\% of native width for HRNet and 1.721\% for matched SegFormer, corresponding to approximately 15.10 and 19.80 native pixels. The absolute difference is 0.409 percentage points of width, or 4.70 pixels. Relative to the SegFormer mean, HRNet reduces error by 23.75\%. SegFormer raw coverage is 100\%, with zero fallback substitutions, so its raw and fallback-assisted means coincide.

HRNet has lower mean error in every seed and on each of the two test tyres (Fig.~\ref{fig:geometry}). This does not mean every individual point improves. The new-tyre example has higher error for both models than the high-mileage example, but two tyres cannot establish a class effect. The result supports HRNet as a candidate point-localisation component within this matched study. It does not establish an architecture-only advantage because supervision and initial representations differ, nor does it demonstrate downstream classification improvement or physical-angle accuracy.
\begin{table}[t]
\caption{Matched fixed-epoch-60 horizontal point error (\% of native width minus one). Each seed uses the same 24 images from two test tyres; coverage is the SegFormer raw point coverage.}
\label{tab:geometry}
\centering
\small
\setlength{\tabcolsep}{4pt}
\begin{tabular}{lrrr}
\toprule
Seed & HRNet & SegFormer & Coverage (\%) \\
\midrule
1 & 1.391 & 1.788 & 100 \\
2 & 1.284 & 1.705 & 100 \\
3 & 1.261 & 1.668 & 100 \\
Mean & 1.312 & 1.721 & 100 \\
\bottomrule
\end{tabular}
\end{table}

\begin{figure*}[t]
\centering
\includegraphics[width=\textwidth]{geometry.pdf}
\caption{Matched boundary comparison. (a) Mean horizontal error for each training seed on the same 24 images. (b) Per-tyre error, retaining all three seed values; horizontal marks are seed means. HRNet is lower in each seed/tyre comparison. These observations are correlated and provide no independent-tyre significance interval.}
\label{fig:geometry}
\end{figure*}

\subsection{Workstation Integration and Observed Video Failures}
Tread Station is a native PySide6/PyTorch application with photo/video loading, portrait-aware playback, asynchronous inference, comparison, geometric overlays, and saved evidence restoration. Existing classifier and region checkpoints remain separate from the optional HRNet/matched-SegFormer geometry components. An inference result is attached to its analysed frame; a newer live preview is shown separately rather than receiving a stale mask. Figure~\ref{fig:station} shows the implemented interface, and Fig.~\ref{fig:qualitative} pairs an operational still example with a preserved video failure.

The saved adapter checks cover strict checkpoint loading, preprocessing, coordinate decoding, repeatability, blank-input flags, temporal resets, and immutable source pixels. A local check on one internal still measured five warm inference calls on an NVIDIA GTX 1650. HRNet's median was 72.60 ms (range 51.41--88.46 ms); matched SegFormer's median was 33.48 ms (33.32--33.84 ms). These timings include preprocessing, inference, extraction, and GPU synchronisation, but exclude model loading and file decoding. Peak allocated memory in the two-component check was approximately 182 MiB. These small operational samples are not a controlled training-hardware comparison, sustained camera throughput, or a guaranteed frame rate.

The supplied videos expose crossed or collapsed HRNet boundaries, missing matched boundaries, and disagreement between models. Crossed pairs withhold widths and connecting geometry while retaining raw proposals for review. No video point ground truth is available, so the passing playback/UI tests establish software behaviour rather than localisation accuracy. Temporal smoothing cannot correct a systematically wrong proposal, and still-image gains cannot be transferred to video by assertion.

The workstation additionally exports the shown analysed frame as PNG and processes an opened video separately at ten output frames per second. The latter runs with a separate CPU engine and frozen display/recipe settings, leaving the GPU available for inspection. The software checks cover frame selection, output duration, portrait dimensions, cancellation, and incomplete-output cleanup; a short real-model portrait excerpt verifies the complete export path. Export frame rate describes the saved playback, not inference speed. This facility supports evidence sharing and does not add an accuracy result.

\begin{figure*}[t]
\centering
\includegraphics[width=.92\textwidth]{video1-learned-workstation.png}
\caption{Saved native workstation view with learned geometry enabled. This operational screenshot predates the additional download controls. The inspected frame, source preview, model observations, and review flags are distinct. Interface completion is not field-accuracy validation.}
\label{fig:station}
\end{figure*}

\begin{figure*}[t]
\centering
\begin{minipage}[t]{.45\textwidth}
\centering
\includegraphics[width=\linewidth,height=6cm,keepaspectratio]{learned-image-check.png}\\
\small (a) Selected internal still-image check
\end{minipage}\hfill
\begin{minipage}[t]{.52\textwidth}
\centering
\includegraphics[width=\linewidth,height=6cm,keepaspectratio]{video1-learned-diagram.png}\\
\small (b) Preserved video-domain failure
\end{minipage}
\caption{Qualitative operational evidence, not a new held-out benchmark. Amber HRNet proposals, cyan matched-SegFormer points, and the image-space centreline are shown on the still. The video explanation retains crossed/missing boundaries and disagreement rather than concealing failure. All widths and slopes are image-relative. The images are reproduced from the saved integration checks without altering predicted coordinates.}
\label{fig:qualitative}
\end{figure*}

\subsection{Alignment Software Is Tested Synthetically}
The target-assisted interface generates distinct ground/wheel boards, creates or loads a camera profile, solves target poses, applies the declared side convention, and records signed outputs with pose residuals. It also accepts an independent reference reading for comparison. Synthetic tests exercise known poses, both wheel sides, missing targets, scaling/orientation, setup gating, and evidence export (Fig.~\ref{fig:alignment}). The supplied videos do not contain this calibrated dual-target setup. No result here measures real wheel-angle error, fixture mounting error, or runout compensation.
